% Oct 26 2026: Figure 7 from output/model/experiments/housing_supply/20261021_free_entry/fig7_fi/fig_m_illustration_data.json
% (notes/model_analysis_rewrite_20261026/scripts/build_fig_m_demand.py).
\subsubsection*{The figures of Section~\ref{sec:m-analysis}.}
\label{app:model-figures}
I compute Figures~\ref{fig:m-risk} and~\ref{fig:m-illustration} from the model at the parameters of
Table~\ref{tab:quant-parameters}. The estimated values are
$\beta_{\text{ann}}=0.9357$, $\psi_0=0.1774$, $\gamma=0.6457$, $\xi=0.2946$,
$\kappa_1=0.0505$, $\kappa_C=0.1285$, $h_P=0.8899$, $\chi=1.1121$,
$\kappa_T=0.0131$ and $\theta_0=0.2435$. The assigned parameters are those of
the same table. In the 2007 steady state at these values, the price
is $P_0=0.7707$, the rebate is $T=0.1879$ per household and there are $N_0=1$
households. Supply \eqref{eq:m-supply} is $H=6.776\,P^{0.63}$. In every solve,
the rebate balances the government budget (relative residual below $10^{-6}$).
The steady-state rent is
\[
 r=\Bigl[1+\tau^p-\frac{1-\delta_H}{R_b}\Bigr]P=0.1698\,P .
\]
The axes of the figures have no numbers, so I report the values at the
marked points below.

\paragraph{Figure~\ref{fig:m-risk}.} This figure comes directly from
\eqref{eq:m-indirect}, without a solve. With $\sigma=2$, \eqref{eq:app-delta}
gives the cost of a first child
\[
 \Delta(y)=\frac{1}{\Theta}\Bigl[\frac{e(1)}{y+T-rh_P}-\frac{1}{y+T}\Bigr],
 \qquad
 \Theta=\alpha_0^{\alpha_0}(1-\alpha_0)^{1-\alpha_0}r^{\alpha_0-1}.
\]
I set $\alpha_0=0.733$, $e(1)=1.35^{0.7}=1.234$ and $h_P=0.890$, and both
curves use the 2007 rebate, $T=0.188$. The solid curve uses the 2007 rent,
$r=0.1308$, so $rh_P=0.116$. The dashed curve uses the rent at the new
steady-state price $P_1=0.408$ of Figure~\ref{fig:m-illustration},
$r=0.0693$, so $rh_P=0.062$. Earnings $y_L=0.274$ and $y_H=0.752$ are
after-tax earnings per period at ages 22--25 in the first and third of the
nine efficiency states, $z=0.103$ and $z=0.284$, and $\bar y$ is their mean.
In the 2007 solve, at least 98\% of renters of that age in these two states
end the period at the borrowing limit. Over the plotted range, renters stay
below the rental cap $\bar h^R=6$, with and without children, at both rents.
The braces mark the extra cost due to risk,
\[
 \tfrac12\bigl[\Delta(y_L)+\Delta(y_H)\bigr]-\Delta(\bar y),
\]
which is 0.245 at the 2007 rent and 0.110 at the new steady-state rent.

\paragraph{Figure~\ref{fig:m-illustration}.} Both panels come from 22
stationary solves at given prices: a solve of the household problem and of
the distribution of households at that price. Ten are at $\psi_0=0.1774$, at
prices from 0.30 to 0.92. The other twelve are at $\psi_1=0.1484$ over the
same range, and three of them are steps of the search for $P_1$. Here $\psi_1$
is the limit $\psi'$ of the decay path \eqref{eq:psi-decay}, with $\rho=0.35$.
At each price, $n(P;\psi)$ is births per entering household, with the bin of
three or more children counted as 3.60 and replacement at $\bar n=2.1$, and
$h(P;\psi)$ is owned plus rented housing per household. The demand and
fertility curves are monotone cubic splines through the solved points, and
supply is $H=6.776\,P^{0.63}$. The marked points are:
\begin{itemize}
 \item $A$, the 2007 steady state: $P_0=0.771$, $H=5.750$, $N_0=1$ and
 $n=\bar n$.
 \item $A'$, preferences $\psi_1$ at the price $P_0$ with $N_0$ households:
 $n=1.654$. In panel (a), demand at $A'$ is $N_0h(P_0;\psi_1)=5.701$, 0.9\%
 below demand at $A$, and on impact the price of the transition in
 Section~\ref{sec:quant-transition} is 1.2\% below $P_0$, so I mark $A'$ at
 $A$.
 \item $B$, the new steady state: $P_1=0.408$, $n=\bar n$ (births within
 $10^{-6}$ of replacement), $H=3.851$, $N_1=H/h(P_1;\psi_1)=0.488$ and
 $T=0.136$.
\end{itemize}
From $A$ to $B$, the price falls by 47\%, the stock by 33\% and the number of
households by 51\%. The arrows show the direction of the adjustment.
